\chapter{Introduction}

\section{Background}
The ability to capture, interpret, and replicate human motion has long been a fundamental challenge in robotics, computer vision, and human-computer interaction. Traditional motion capture systems rely on expensive marker-based setups, wearable inertial measurement units (IMUs), or specialized depth sensors such as the Microsoft Kinect and Intel RealSense, all of which impose significant hardware constraints and deployment costs. Recent advances in deep learning-based pose estimation, particularly lightweight frameworks such as MediaPipe developed by Google, have enabled accurate real-time extraction of human joint coordinates from standard RGB video without specialized hardware.
 
Simultaneously, physics-based simulation platforms such as PyBullet have matured into robust environments for replicating articulated body motion, providing physically accurate rendering of joint movement and a foundation for generating structured motion datasets. The convergence of these two technologies — vision-based pose estimation and physics-based simulation — presents an opportunity to build affordable, accessible motion capture systems that require only standard cameras.
 
In the context of industrial robotics and assistive automation, large-scale datasets of human arm movement are essential for training robotic manipulators to perform tasks mimicking natural human motion. However, the cost and complexity of existing data collection pipelines remain a barrier to widespread dataset generation. A vision-based, multi-camera system capable of capturing full arm-to-fingertip motion in real-time and replicating it in simulation addresses this barrier directly.

\section{Motivation}
Robotic arms capable of replicating human arm motion have significant potential in industrial automation, assistive robotics, teleoperation, and rehabilitation engineering. A core requirement for training such systems is access to large, diverse datasets of human arm trajectories annotated with joint coordinates and angles. Existing approaches to generating such datasets either rely on expensive hardware (RGB-D sensors, IMU suits) or produce offline, non-real-time outputs that are difficult to scale.
 
The motivation for MotionBridge arises from the absence of a low-cost, real-time, vision-only pipeline that combines multi-camera capture, neural network-based pose estimation, three-dimensional reconstruction, and physics simulation into a single integrated system. By using only standard USB webcams and open-source software, the proposed system can be deployed in resource-constrained environments such as academic laboratories, small research groups, and developing-world institutions where expensive motion capture infrastructure is unavailable.
 
Furthermore, no existing reviewed work captures both full arm joint trajectories and detailed finger joint coordinates (21 points per hand) from standard RGB cameras and simultaneously replicates arm movement in simulation. MotionBridge addresses this gap, producing a richer dataset than currently available vision-based approaches.

\section{Problem Statement}
Existing human arm motion capture systems for robotic dataset generation suffer from one or more of the following limitations: reliance on expensive depth-sensing hardware, dependence on wearable sensors, absence of real-time simulation replication, operation on monocular camera inputs that cannot resolve depth accurately, or omission of finger-level joint data. These constraints prevent scalable and affordable motion dataset generation for robotic arm training.
 
This project addresses the problem of building a low-cost, real-time system that captures three-dimensional human arm and finger joint motion using only standard RGB cameras, replicates the captured motion in a physics-based simulation environment, and produces a structured dataset of joint coordinates and angles for downstream robotic arm training — without requiring specialized depth sensors, wearable hardware, or GPU-accelerated inference. 

\section{Objectives}



\section{Project Scope}
\subsection{Capabilities}
\begin{itemize}
    \item Real-time capture of human arm movement (shoulder, elbow, wrist) using three standard USB webcams.
    \item Detection of 21 finger joint landmarks per hand using MediaPipe Hands in addition to full arm keypoints.
    \item Three-dimensional joint coordinate reconstruction via confidence-weighted triangulation using calibrated multi-camera geometry.
    \item Real-time replication of arm movement in PyBullet physics simulation using Inverse Kinematics.
    \item Generation of a structured dataset containing per-frame 3D joint coordinates, joint angles, and timestamps.
    \item Occlusion robustness through multi-camera redundancy when one camera's view is blocked, remaining cameras compensate.
\end{itemize}

\subsection{Limitations}
\begin{itemize}
    \item Finger joint movement is captured and recorded in the dataset but is not replicated in PyBullet simulation, as standard arm URDF models do not include articulated finger structures.
    \item The system is limited to single-person capture; multi-person simultaneous tracking is outside scope.
    \item Accuracy depends on initial camera calibration quality; dynamic camera repositioning requires recalibration.
    \item The system does not include robot arm training or policy learning; it generates the dataset only.
    \item Performance may degrade under extreme lighting conditions or significant self-occlusion not resolvable by three cameras.
    \item The system does not support full-body motion capture only the arm from shoulder to fingertip.
\end{itemize}


\section{Problem Applcations}
The dataset and simulation output produced by MotionBridge have direct applications in the following domains:
 
\begin{itemize}
    \item \textbf{Industrial Automation:} Providing training data for robotic arms performing labour-intensive factory tasks such as assembly, sorting, and load-carrying that mimic human arm motion.
    \item \textbf{Assistive Robotics:} Generating motion trajectories for robotic prosthetics and assistive devices that replicate natural human arm and hand movement.
    \item \textbf{Teleoperation:} Serving as a reference system for mapping human operator arm motion to remote robotic manipulators.
    \item \textbf{Rehabilitation Engineering:} Capturing and analysing patient arm movement for physiotherapy assessment and progress monitoring.
    \item \textbf{Imitation Learning Research:} Supplying labelled demonstration datasets for behavioral cloning and inverse reinforcement learning pipelines.
    \item \textbf{Simulation-Based Training:} Using the PyBullet output as a reference motion environment for reinforcement learning agents.
\end{itemize} 
